\begin{proof}
Here $T\mathbf1=\frac12\mathbf1$, so termination is geometric.
The two-row reachability matrix with rows $\alpha,\alpha A_0$
and the two-column observability matrix with columns $\tau,A_0\tau$
both have determinant $-1/40$, proving rank two.
For this example write $f_t(v)=(t,1-t)A_v\tau$; its first coordinate
updates by
\[
 F_0(t)=\frac{1+t}{2+t},\qquad F_1(t)=\frac1{3-t}.
\]
On $[0,1]$ their images are $[1/2,2/3]$ and $[1/3,1/2]$.
With $r=(\sqrt5-1)/2$, both maps preserve $(1-r,r)$, and every
finite-prefix coordinate lies in that interval. The words $0^n$
and $1^n$ approach its two endpoints. For any finite dictionary let
$m<r$ be its maximum coordinate. Then
\[
 F_0(m)-m=\frac{1-m-m^2}{2+m}>0.
\]
The next residual is outside the dictionary's convex hull. The
coordinate map is affine and injective since
$f_t(0)=1/10+t/20$. Nonnegative representations of normalized laws
have coefficient sum one, so convex and conic closure give the same
obstruction after normalizing each nonzero shift.
\end{proof}
